\documentclass{amsart}
\usepackage{amsmath,amssymb,amsthm,hyperref}
\newtheorem{theorem}{Theorem}
\newtheorem{lemma}{Lemma}
\newtheorem{proposition}{Proposition}
\theoremstyle{definition}
\newtheorem{definition}{Definition}
\theoremstyle{remark}
\newtheorem{remark}{Remark}
\begin{document}

\title{The moduli space of left-invariant statistical connections on the Gaussian Lie group}
\maketitle

\section{Setup and statement of the result}

Let $G=\{(\mu,\sigma):\mu\in\mathbb{R},\ \sigma>0\}$ with product
$(\mu_1,\sigma_1)\cdot(\mu_2,\sigma_2)=(\mu_1+\sigma_1\mu_2,\ \sigma_1\sigma_2)$ and Fisher metric
$g=\sigma^{-2}\,(d\mu\otimes d\mu+2\,d\sigma\otimes d\sigma)$.
As a manifold $G\cong\mathbb{R}\times\mathbb{R}_{>0}$ is simply connected.

\begin{theorem}\label{thm:main}
There is a bijection
\[
\mathcal{C}(G,g)\ \xrightarrow{\ \cong\ }\ \mathbb{R}^4,\qquad \nabla\longmapsto (a,b,d,e),
\]
under which $\nabla$ corresponds to the totally symmetric cubic form with components
$(C_{111},C_{112},C_{122},C_{222})=(a,b,d,e)$ in the orthogonal left-invariant frame of
Section~\ref{sec:frame}. The group of $g$-isometric Lie-group automorphisms of $G$ is
\[
\operatorname{Aut}_{\mathrm{isom}}(G,g)=\{\mathrm{id},\ F\},\qquad F(\mu,\sigma)=(-\mu,\sigma),
\]
a group isomorphic to $\mathbb{Z}/2\mathbb{Z}$, and its pullback action on $\mathcal{C}(G,g)\cong\mathbb{R}^4$ is
\begin{equation}\label{eq:action}
(a,b,d,e)\ \longmapsto\ (-a,\,b,\,-d,\,e).
\end{equation}
Consequently the moduli space is
\[
\mathcal{M}(G,g)=\mathbb{R}^4/(\mathbb{Z}/2\mathbb{Z})
\ \cong\ \mathbb{R}^2_{(b,e)}\times\bigl(\mathbb{R}^2_{(a,d)}/\{\pm 1\}\bigr).
\]
Explicitly, using the invariants $b,\ e,\ s:=a^2,\ t:=d^2,\ r:=ad$, one has the algebraic model
\[
\mathcal{M}(G,g)\ \cong\ \bigl\{(b,e,s,t,r)\in\mathbb{R}^5:\ s\ge 0,\ t\ge 0,\ st=r^2\bigr\}.
\]
The latter is the product of the plane $\mathbb{R}^2_{(b,e)}$ with a two-dimensional cone (the quotient of the
plane by the antipodal involution); it is a real $4$-dimensional orbifold with a single conical singular
stratum $\{a=d=0\}\cong\mathbb{R}^2_{(b,e)}$.
\end{theorem}

The proof occupies Sections~\ref{sec:frame}--\ref{sec:quotient}.

\section{A left-invariant orthogonal frame and the Lie algebra}\label{sec:frame}

\begin{lemma}\label{lem:frame}
The vector fields
\[
e_1=\sigma\,\partial_\mu,\qquad e_2=\sigma\,\partial_\sigma
\]
are left-invariant, satisfy
\begin{equation}\label{eq:metric}
g(e_1,e_1)=1,\qquad g(e_2,e_2)=2,\qquad g(e_1,e_2)=0,
\end{equation}
and
\begin{equation}\label{eq:bracket}
[e_1,e_2]=-e_1.
\end{equation}
\end{lemma}

\begin{proof}
For $(a,b)\in G$ the left translation is $L_{(a,b)}(\mu,\sigma)=(a+b\mu,\,b\sigma)$, with constant
differential $d L_{(a,b)}=\mathrm{diag}(b,b)$ in the coordinate basis $(\partial_\mu,\partial_\sigma)$.
A left-invariant field $X$ is determined by its value $X_e$ at the identity $e=(0,1)$ via
$X_{(\mu,\sigma)}=dL_{(\mu,\sigma)}(X_e)$. Taking $X_e=\partial_\mu$ gives
$dL_{(\mu,\sigma)}\partial_\mu=\sigma\,\partial_\mu=e_1$, and taking $X_e=\partial_\sigma$ gives
$dL_{(\mu,\sigma)}\partial_\sigma=\sigma\,\partial_\sigma=e_2$. Hence $e_1,e_2$ are left-invariant.

For \eqref{eq:metric}: $g(e_1,e_1)=\sigma^2\cdot\sigma^{-2}=1$, $g(e_2,e_2)=\sigma^2\cdot 2\sigma^{-2}=2$,
and $g(e_1,e_2)=0$ since $g$ has no cross term. (These constants confirm $g$ is left-invariant:
its components in the left-invariant frame are constant.)

For \eqref{eq:bracket}: for any smooth $f$,
\[
[e_1,e_2]f=e_1(e_2 f)-e_2(e_1 f)
=\sigma\partial_\mu(\sigma f_\sigma)-\sigma\partial_\sigma(\sigma f_\mu)
=\sigma^2 f_{\mu\sigma}-\sigma f_\mu-\sigma^2 f_{\sigma\mu}=-\sigma f_\mu=-e_1 f. \qedhere
\]
\end{proof}

The Lie algebra $\mathfrak{g}=\operatorname{span}(e_1,e_2)$ with \eqref{eq:bracket} is the non-abelian
two-dimensional Lie algebra $\mathfrak{aff}(1)$. The metric \eqref{eq:metric} is encoded by the constant
symmetric matrix
\begin{equation}\label{eq:Gmat}
\mathbf{G}=\begin{pmatrix}1&0\\0&2\end{pmatrix},\qquad \mathbf{G}_{ij}=g(e_i,e_j).
\end{equation}

\section{Left-invariant connections in the frame}\label{sec:conn}

A left-invariant affine connection $\nabla$ is determined by constants
$\Gamma^{k}_{ij}\in\mathbb{R}$ via
\[
\nabla_{e_i}e_j=\sum_{k}\Gamma^{k}_{ij}\,e_k,\qquad i,j\in\{1,2\},
\]
and conversely any choice of these eight constants defines a left-invariant connection. (This is the standard
correspondence between left-invariant connections and bilinear maps $\mathfrak g\times\mathfrak g\to\mathfrak g$:
left invariance of $\nabla$ is equivalent to the $\Gamma^k_{ij}$ being constant in the global left-invariant
frame $(e_1,e_2)$.)

\subsection*{Torsion-freeness}
Let $c^{k}_{ij}$ denote the structure constants, $[e_i,e_j]=\sum_k c^{k}_{ij}e_k$; by \eqref{eq:bracket} the only
nonzero ones are $c^{1}_{12}=-1$, $c^{1}_{21}=1$. The torsion of $\nabla$ is
$T(e_i,e_j)=\nabla_{e_i}e_j-\nabla_{e_j}e_i-[e_i,e_j]$, with components
\begin{equation}\label{eq:torsion}
T^{k}_{ij}=\Gamma^{k}_{ij}-\Gamma^{k}_{ji}-c^{k}_{ij}.
\end{equation}

\subsection*{Compatibility and the dual connection}
Because $\mathbf{G}_{ij}=g(e_i,e_j)$ is constant, we have $e_i\,g(e_j,e_k)=0$ for all $i,j,k$. Hence for a
left-invariant connection $\nabla$ and a left-invariant candidate dual $\nabla^{*}$ with constants
$\Gamma^{*\,k}_{ij}$, the duality identity
$X\,g(Y,Z)=g(\nabla_X Y,Z)+g(Y,\nabla^{*}_X Z)$ evaluated on the frame ($X=e_i,Y=e_j,Z=e_k$) reads, for all
$i,j,k$,
\begin{equation}\label{eq:compat}
0=\sum_{l}\mathbf{G}_{lk}\,\Gamma^{l}_{ij}+\sum_{l}\mathbf{G}_{jl}\,\Gamma^{*\,l}_{ik}.
\end{equation}
For fixed $\nabla$, \eqref{eq:compat} \emph{uniquely} determines the constants $\Gamma^{*\,l}_{ik}$ (solve using
$\mathbf G$ invertible), so every left-invariant $\nabla$ has a unique left-invariant dual.
We record that admitting \emph{some} affine dual is equivalent to admitting this left-invariant one: if
$\widetilde\nabla^{*}$ is any affine connection satisfying the duality identity, then both
$\widetilde\nabla^{*}$ and the left-invariant solution $\nabla^{*}$ above satisfy
$g(Y,\widetilde\nabla^{*}_XZ)=X g(Y,Z)-g(\nabla_XY,Z)=g(Y,\nabla^{*}_XZ)$ for all $X,Y,Z$, hence
$\widetilde\nabla^{*}=\nabla^{*}$ since $g$ is nondegenerate. Thus seeking the dual among left-invariant
connections loses no generality, and ``$\nabla$ admits a torsion-free dual'' is unambiguous.

\begin{proposition}\label{prop:dim}
The set $\mathcal{C}(G,g)$ of left-invariant torsion-free connections $\nabla$ admitting a torsion-free dual
$\nabla^{*}$ is in bijection with the space of totally symmetric trilinear forms
$C$ on $\mathfrak g$, via
\begin{equation}\label{eq:cubic}
C_{kij}:=g\!\left(\,\nabla_{e_i}e_j-\nabla^{g}_{e_i}e_j,\,e_k\right)
=\sum_{l}\mathbf{G}_{kl}\bigl(\Gamma^{l}_{ij}-\Gamma^{g\,l}_{ij}\bigr),
\end{equation}
where $\nabla^{g}$ is the Levi--Civita connection of $g$. In particular $\mathcal{C}(G,g)$ is a real vector space
of dimension $\binom{2+3-1}{3}=4$, with coordinates
$(a,b,d,e)=(C_{111},C_{112},C_{122},C_{222})$.
\end{proposition}

\begin{proof}
\emph{Step 1: the Levi--Civita connection.}
Since $\mathbf G$ is constant on the frame, the Koszul formula reduces (the three derivative terms vanish) to
\[
2\,g(\nabla^{g}_{e_i}e_j,\,e_m)=g([e_i,e_j],e_m)-g([e_j,e_m],e_i)+g([e_m,e_i],e_j).
\]
Using \eqref{eq:bracket} and \eqref{eq:Gmat} one computes
\begin{equation}\label{eq:LC}
\nabla^{g}_{e_1}e_1=\tfrac12\,e_2,\quad
\nabla^{g}_{e_1}e_2=-\,e_1,\quad
\nabla^{g}_{e_2}e_1=0,\quad
\nabla^{g}_{e_2}e_2=0.
\end{equation}
For instance, $2g(\nabla^g_{e_1}e_1,e_2)=-g([e_1,e_2],e_1)+g([e_2,e_1],e_1)=2\,g(e_1,e_1)=2$ gives the
$e_2$-component $\tfrac12$, while $2g(\nabla^g_{e_1}e_1,e_1)=g([e_1,e_1],e_1)=0$ gives the $e_1$-component $0$;
the other three covariant derivatives are obtained the same way. The connection $\nabla^g$ is torsion-free and
metric.

\emph{Step 2: reduction to the difference tensor and its dual.}
Write any left-invariant connection as $\nabla=\nabla^{g}+A$, where
$A(e_i,e_j):=\nabla_{e_i}e_j-\nabla^g_{e_i}e_j$ is a left-invariant $(1,2)$-tensor with constant components
$A^{l}_{ij}=\Gamma^{l}_{ij}-\Gamma^{g\,l}_{ij}$. Lowering the contravariant index, set
$C_{kij}:=\sum_l \mathbf G_{kl}A^{l}_{ij}$, in agreement with \eqref{eq:cubic}.
Because $\nabla^g$ is metric, $e_i\,g(e_j,e_k)=g(\nabla^g_{e_i}e_j,e_k)+g(e_j,\nabla^g_{e_i}e_k)$. Subtracting
this identity from the duality identity \eqref{eq:compat} shows that the dual of $\nabla=\nabla^g+A$ is
$\nabla^{*}=\nabla^{g}+A^{*}$, where the constants $A^{*\,l}_{ij}$ satisfy
\begin{equation}\label{eq:compatA}
\sum_l\mathbf G_{lk}\,A^{l}_{ij}+\sum_l\mathbf G_{jl}\,A^{*\,l}_{ik}=0\qquad\text{for all }i,j,k.
\end{equation}
Define the lowered tensor $C^{*}_{kij}:=\sum_l\mathbf G_{kl}A^{*\,l}_{ij}$. The first sum in
\eqref{eq:compatA} equals $C_{kij}$, and the second equals $C^{*}_{jik}$ (it lowers the index $l$ against the
\emph{first} free slot $j$ of $A^*$). Hence \eqref{eq:compatA} is exactly
\begin{equation}\label{eq:CR}
C_{kij}+C^{*}_{jik}=0\qquad\text{for all }i,j,k.
\end{equation}
Relabelling the free indices in \eqref{eq:CR} (put $A=j$, $B=i$, $C=k$, so $C^*_{jik}=C^*_{ABC}$ and
$C_{kij}=C_{CBA}$) gives the equivalent closed form of the compatibility relation:
\begin{equation}\label{eq:dualA}
C^{*}_{ABC}=-\,C_{CBA}\qquad\text{(reverse the three slots and negate).}
\end{equation}
Equation \eqref{eq:dualA} holds for an \emph{arbitrary} left-invariant $\nabla$ and its unique dual; no symmetry
of $C$ has yet been used.

\emph{Step 3: the two torsion-free conditions force total symmetry of $C$.}
By \eqref{eq:torsion} and torsion-freeness of $\nabla^g$, the connection $\nabla=\nabla^g+A$ is torsion-free iff
$A^{l}_{ij}=A^{l}_{ji}$ for all $i,j,l$, i.e. (lowering with the invertible $\mathbf G$) iff
\begin{equation}\label{eq:P}
C_{xyz}=C_{xzy}\qquad\text{(symmetry in slots $2,3$).}
\end{equation}
Likewise $\nabla^{*}$ is torsion-free iff $C^{*}_{xyz}=C^{*}_{xzy}$. Substituting \eqref{eq:dualA}, this reads
$-C_{zyx}=-C_{yzx}$, i.e.
\begin{equation}\label{eq:Q}
C_{zyx}=C_{yzx}\qquad\text{(symmetry in slots $1,2$).}
\end{equation}
The transposition of slots $(2,3)$ in \eqref{eq:P} and the transposition of slots $(1,2)$ in \eqref{eq:Q}
generate the full symmetric group $S_3$ acting on the three slots. Therefore $C$ is invariant under all
permutations of its indices, i.e. $C$ is totally symmetric. (Explicitly: \eqref{eq:P} and \eqref{eq:Q} together
imply $C_{112}=C_{121}=C_{211}$ and $C_{122}=C_{212}=C_{221}$, leaving the four independent components
$C_{111},C_{112},C_{122},C_{222}$.)

\emph{Step 4: converse and dimension count.}
Conversely, let $C$ be any totally symmetric trilinear form on $\mathfrak g$. Define $\nabla=\nabla^g+A$ by
$A^{l}_{ij}=\sum_k\mathbf G^{lk}C_{kij}$ (raising the first index), and $\nabla^{*}=\nabla^g+A^{*}$ by
$A^{*\,l}_{ij}=\sum_k\mathbf G^{lk}C^{*}_{kij}$ with $C^{*}_{ABC}:=-C_{CBA}$ as in \eqref{eq:dualA}. Then:
\begin{itemize}
\item Both connections are left-invariant (constant components in the frame).
\item $\nabla$ is torsion-free because $C$ is symmetric in slots $2,3$ (total symmetry); $\nabla^{*}$ is
torsion-free because $C^{*}_{xyz}=-C_{zyx}=-C_{zxy}=C^{*}_{xzy}$, using total symmetry of $C$.
\item $\nabla$ and $\nabla^{*}$ are mutually dual: with $C^{*}_{ABC}=-C_{CBA}$, the compatibility relation
\eqref{eq:CR} reads $C_{kij}+C^{*}_{jik}=C_{kij}-C_{kij}=0$ (here $C^{*}_{jik}=-C_{kij}$ by \eqref{eq:dualA}),
which is exactly \eqref{eq:compatA}; tracing back, $\nabla,\nabla^{*}$ satisfy the duality identity on the frame,
hence (by left invariance and tensoriality, see the discussion before this proposition) for all vector fields.
\end{itemize}
Therefore $\nabla\in\mathcal C(G,g)$. The assignment $\nabla\mapsto C$ of Step 2 and the assignment
$C\mapsto\nabla$ of Step 4 are mutually inverse linear maps, so $\mathcal C(G,g)$ is linearly isomorphic to the
space of totally symmetric trilinear forms on the $2$-dimensional space $\mathfrak g$. That space has dimension
$\binom{2+3-1}{3}=4$, with coordinates $(a,b,d,e)=(C_{111},C_{112},C_{122},C_{222})$.
\end{proof}

\begin{remark}
Equivalently, the dual pairs are $\nabla^{(\pm)}=\nabla^g\mp\tfrac12 C^{\sharp}$ with $C$ the Amari--Chentsov
cubic form (totally symmetric), where $C^{\sharp}$ raises the first index of $C$. This recovers the classical
$\alpha$-family $C\mapsto\alpha C$ as the line through any fixed nonzero $C$.
\end{remark}

\section{The isometric automorphism group}\label{sec:aut}

\begin{lemma}\label{lem:aut}
Let $\operatorname{Aut}(G)$ be the group of Lie-group automorphisms of $G$. The differential at $e$ gives an
isomorphism $\operatorname{Aut}(G)\cong\operatorname{Aut}(\mathfrak g)$, and
\[
\operatorname{Aut}(\mathfrak g)=\Bigl\{\ \phi:\ e_1\mapsto p\,e_1,\ \ e_2\mapsto q\,e_1+e_2\ \Bigm|\ p\in\mathbb R^\times,\ q\in\mathbb R\ \Bigr\}.
\]
Among these, the $g$-isometric ones (those whose differential preserves $\mathbf G$) are exactly
$\phi_{\pm}:e_1\mapsto\pm e_1,\ e_2\mapsto e_2$. The corresponding group automorphisms are $\mathrm{id}$ and
$F(\mu,\sigma)=(-\mu,\sigma)$, and
\[
\operatorname{Aut}_{\mathrm{isom}}(G,g)=\{\mathrm{id},F\}\cong\mathbb{Z}/2\mathbb{Z}.
\]
\end{lemma}

\begin{proof}
\emph{Group vs.\ algebra automorphisms.} $G$ is connected and simply connected (diffeomorphic to
$\mathbb R\times\mathbb R_{>0}$). For a connected, simply connected Lie group the differential at the identity
$F\mapsto d F_e$ is a group isomorphism $\operatorname{Aut}(G)\xrightarrow{\cong}\operatorname{Aut}(\mathfrak g)$:
every Lie-algebra automorphism integrates to a unique Lie-group automorphism. Hence we may compute on
$\mathfrak g$.

\emph{Computing $\operatorname{Aut}(\mathfrak g)$.} The derived algebra $[\mathfrak g,\mathfrak g]$ is spanned by
$[e_1,e_2]=-e_1$, i.e.\ equals $\mathbb R e_1$. Any $\phi\in\operatorname{Aut}(\mathfrak g)$ preserves the derived
algebra, so $\phi(e_1)=p\,e_1$ with $p\neq0$, and $\phi(e_2)=q\,e_1+r\,e_2$ with $r\neq0$ (otherwise $\phi$ is not
invertible). Then
\[
\phi([e_1,e_2])=\phi(-e_1)=-p\,e_1,\qquad
[\phi(e_1),\phi(e_2)]=[p e_1,\,q e_1+r e_2]=pr\,[e_1,e_2]=-pr\,e_1.
\]
Equality forces $pr=p$, hence $r=1$. This gives precisely the stated family.

\emph{Isometric condition.} A Lie-group automorphism $F$ satisfies $F\circ L_h=L_{F(h)}\circ F$ for all $h\in G$;
together with left-invariance of $g$ this implies $F^{*}g$ is again left-invariant, and $F$ is a $g$-isometry iff
its differential at $e$ preserves the inner product $\mathbf G$ on $\mathfrak g$. Writing
$\phi(e_1)=p\,e_1,\ \phi(e_2)=q\,e_1+e_2$ as the matrix
$M=\left(\begin{smallmatrix}p&q\\0&1\end{smallmatrix}\right)$ in the basis $(e_1,e_2)$, the condition
$M^{\mathsf T}\mathbf G\,M=\mathbf G$ reads
\[
\begin{pmatrix}p^2 & pq\\ pq & q^2+2\end{pmatrix}=\begin{pmatrix}1&0\\0&2\end{pmatrix},
\]
giving $p^2=1$, $pq=0$, $q^2+2=2$, hence $q=0$ and $p=\pm1$. Thus the isometric algebra automorphisms are
$\phi_{\pm}$.

\emph{Integration.} $\phi_+=\mathrm{id}$ integrates to $\mathrm{id}$. For $\phi_-:e_1\mapsto-e_1,\ e_2\mapsto e_2$
the map $F(\mu,\sigma)=(-\mu,\sigma)$ is a Lie-group automorphism:
\[
F\bigl((\mu_1,\sigma_1)\cdot(\mu_2,\sigma_2)\bigr)=F(\mu_1+\sigma_1\mu_2,\ \sigma_1\sigma_2)
=(-\mu_1-\sigma_1\mu_2,\ \sigma_1\sigma_2),
\]
\[
F(\mu_1,\sigma_1)\cdot F(\mu_2,\sigma_2)=(-\mu_1,\sigma_1)\cdot(-\mu_2,\sigma_2)
=(-\mu_1+\sigma_1(-\mu_2),\ \sigma_1\sigma_2)=(-\mu_1-\sigma_1\mu_2,\ \sigma_1\sigma_2),
\]
which agree. Its differential at $e=(0,1)$ sends $\partial_\mu\mapsto-\partial_\mu$,
$\partial_\sigma\mapsto\partial_\sigma$, i.e.\ $e_1\mapsto -e_1,\ e_2\mapsto e_2$, so $dF_e=\phi_-$. Finally $F$
is a $g$-isometry because $g$ is unchanged by $\mu\mapsto-\mu$. Since $F^2=\mathrm{id}$,
$\operatorname{Aut}_{\mathrm{isom}}(G,g)=\{\mathrm{id},F\}\cong\mathbb Z/2\mathbb Z$.
\end{proof}

\section{The action on connections and the quotient}\label{sec:quotient}

\begin{lemma}\label{lem:actC}
The pullback action of $F$ on $\mathcal C(G,g)\cong\mathbb R^4$ is \eqref{eq:action}:
$(a,b,d,e)\mapsto(-a,b,-d,e)$.
\end{lemma}

\begin{proof}
Let $\nabla\in\mathcal C(G,g)$ with cubic form $C$. The pullback connection $F^{*}\nabla$, defined by
$(F^{*}\nabla)_XY=(F^{-1})_{*}\bigl(\nabla_{F_{*}X}F_{*}Y\bigr)$, is again left-invariant (because $F$ is a group
automorphism, it conjugates left translations to left translations, so it preserves left-invariant tensors) and
torsion-free (pullback by a diffeomorphism intertwines torsion tensors). Its dual is $F^{*}(\nabla^{*})$: pulling
back the duality identity $X g(Y,Z)=g(\nabla_XY,Z)+g(Y,\nabla^{*}_XZ)$ by $F$ and using $F^{*}g=g$ (Lemma
\ref{lem:aut}) shows $F^{*}\nabla$ and $F^{*}\nabla^{*}$ satisfy the duality identity, and both are torsion-free.
Hence $F^{*}\nabla\in\mathcal C(G,g)$, and its cubic form is the pulled-back form $F^{*}C$, where
$(F^{*}C)(X,Y,Z)=C(F_{*}X,F_{*}Y,F_{*}Z)$.

Since $C$ is a left-invariant $(0,3)$-tensor and $F$ is a Lie-group automorphism with constant differential
$\phi:=dF_e$ on the left-invariant frame, on the frame
\[
(F^{*}C)_{ijk}=\sum_{\alpha,\beta,\gamma}\phi^{\alpha}{}_{i}\,\phi^{\beta}{}_{j}\,\phi^{\gamma}{}_{k}\,C_{\alpha\beta\gamma},
\]
where $\phi=\mathrm{diag}(-1,1)$ in the basis $(e_1,e_2)$, i.e.\ $\phi^1{}_1=-1$, $\phi^2{}_2=1$, off-diagonal
entries $0$. Thus $(F^{*}C)_{ijk}=(-1)^{\#\{\text{indices equal to }1\}}\,C_{ijk}$. Applied to the four
independent components:
\[
a=C_{111}\mapsto(-1)^3a=-a,\quad
b=C_{112}\mapsto(-1)^2b=b,\quad
d=C_{122}\mapsto(-1)^1d=-d,\quad
e=C_{222}\mapsto e,
\]
where total symmetry of $C$ ensures, e.g., $C_{112}=C_{121}=C_{211}$ all transform identically. This is exactly
\eqref{eq:action}.
\end{proof}

\begin{proof}[Proof of Theorem~\ref{thm:main}]
By Proposition~\ref{prop:dim}, $\mathcal C(G,g)\cong\mathbb R^4$ via $\nabla\mapsto(a,b,d,e)$. By
Lemma~\ref{lem:aut}, the isometric automorphism group is $\mathbb Z/2\mathbb Z=\{\mathrm{id},F\}$, and by
Lemma~\ref{lem:actC} the nontrivial element acts by $\iota:(a,b,d,e)\mapsto(-a,b,-d,e)$. Therefore
\[
\mathcal M(G,g)=\mathcal C(G,g)/\!\sim\ =\ \mathbb R^4/\langle\iota\rangle.
\]
The involution $\iota$ fixes the coordinates $b,e$ and acts as the antipodal map $-\mathrm{id}$ on the plane
$\mathbb R^2_{(a,d)}$. Since the two factors $\mathbb R^2_{(b,e)}$ and $\mathbb R^2_{(a,d)}$ are
$\iota$-invariant and $\iota$ is trivial on the first,
\[
\mathbb R^4/\langle\iota\rangle\ \cong\ \mathbb R^2_{(b,e)}\times\bigl(\mathbb R^2_{(a,d)}/\{\pm 1\}\bigr).
\]

\emph{Explicit algebraic model.} The ring of $\{\pm1\}$-invariant polynomials on $\mathbb R^2_{(a,d)}$ (for the
action $(a,d)\mapsto(-a,-d)$) is generated by $s=a^2$, $t=d^2$, $r=ad$, subject to the single relation $st=r^2$
and the inequalities $s,t\ge0$. The map
\[
\Phi:\ \mathbb R^2_{(a,d)}\longrightarrow\bigl\{(s,t,r)\in\mathbb R^3:\ s,t\ge0,\ st=r^2\bigr\},\qquad
\Phi(a,d)=(a^2,d^2,ad),
\]
is surjective: given $(s,t,r)$ with $s,t\ge0$ and $st=r^2$, choose $a=\sqrt s$ and $d=\operatorname{sgn}(r)\sqrt t$
if $r\neq0$, and $d=\sqrt t$ if $r=0$; then $a^2=s$, $d^2=t$, $ad=\operatorname{sgn}(r)\sqrt{st}=r$. Moreover
$\Phi(a,d)=\Phi(a',d')$ iff $a'^2=a^2$, $d'^2=d^2$, $a'd'=ad$, which (checking signs) holds iff
$(a',d')=\pm(a,d)$. Hence $\Phi$ induces a bijection $\mathbb R^2_{(a,d)}/\{\pm1\}\cong\{(s,t,r):s,t\ge0,\,st=r^2\}$.
Adjoining the free coordinates $b,e$ gives the bijection
\[
\mathcal M(G,g)\ \cong\ \bigl\{(b,e,s,t,r)\in\mathbb R^5:\ s\ge0,\ t\ge0,\ st=r^2\bigr\},\qquad
(a,b,d,e)\mapsto(b,e,a^2,d^2,ad),
\]
which is well defined and injective on $\sim$-classes by the above.

\emph{Geometric description.} The locus $\{a=d=0\}$, i.e.\ $s=t=r=0$, is the fixed-point set of $\iota$; it is a
copy of $\mathbb R^2_{(b,e)}$ and forms the conical singular stratum of the quotient. On the complement
$\{(a,d)\neq(0,0)\}$ the group $\{\pm1\}$ acts freely and properly, so
$\bigl(\mathbb R^2_{(a,d)}\setminus\{0\}\bigr)/\{\pm1\}$ is a smooth $2$-manifold; together with the cone point
this realizes $\mathbb R^2/\{\pm1\}$ as the standard $2$-dimensional cone (the cone over $\mathbb{RP}^1$). Hence
$\mathcal M(G,g)$ is a real $4$-dimensional orbifold: the product of the plane $\mathbb R^2_{(b,e)}$ with this
cone, smooth except along the $\mathbb Z/2$ conical stratum $\{a=d=0\}\cong\mathbb R^2_{(b,e)}$. This is the
required explicit parametrization of the moduli space.
\end{proof}

\end{document}
